\chapter{Paralysis}

\section*{Definition}
Paralysis refers to the loss of voluntary motor function in a body region due to disruption of the motor pathway from the cerebral cortex to the muscle. Hemiparesis is weakness affecting one side of the body; hemiplegia indicates complete loss of movement on one side. Monoplegia involves a single limb, paraplegia affects both lower extremities, and quadriplegia involves all four limbs. The pattern of weakness localizes the lesion along the neuraxis.

\section*{What Happens in Your Body}
Upper motor neuron (UMN) lesions involve the corticospinal tract from motor cortex through the internal capsule, brainstem, and spinal cord, producing spastic weakness with hyperreflexia, Babinski sign, and increased tone. Lower motor neuron (LMN) lesions involve the anterior horn cell, nerve root, or peripheral nerve, causing flaccid weakness with hyporeflexia, fasciculations, and muscle atrophy. The corticospinal tract decussates at the caudal medulla, so lesions above this decussation produce contralateral weakness, while lesions below produce ipsilateral weakness.

\section*{Causes}
\begin{itemize}
  \item Cerebrovascular (most common for acute hemiparesis): Ischemic stroke (large vessel atherosclerosis, cardioembolism, small vessel lacunar), Intracerebral hemorrhage, Subarachnoid hemorrhage; transient ischemic attack (TIA) causes temporary deficits
  \item Structural lesions: Brain tumor (primary or metastatic), Brain abscess, Subdural or epidural hematoma, Spinal cord compression (herniated disc, spinal stenosis, epidural abscess, vertebral metastasis, trauma)
  \item Demyelinating disease: Multiple sclerosis (relapsing-remitting episodes), Transverse myelitis, Acute disseminated encephalomyelitis (ADEM)
  \item Neuromuscular causes: Guillain-Barr\'e syndrome (ascending paralysis), Myasthenia gravis (fluctuating, fatiguable weakness), Amyotrophic lateral sclerosis (ALS, combined UMN/LMN signs), Spinal muscular atrophy, Periodic paralysis (hypokalemic, hyperkalemic), Botulism, Tick paralysis
  \item Infectious: Polio (now rare), West Nile virus, Enterovirus D68 (acute flaccid myelitis), Post-infectious polyneuropathy, HIV-associated myelopathy
  \item Traumatic: Spinal cord injury, Traumatic brain injury, Brachial or lumbosacral plexopathy, Peripheral nerve transection
  \item Other: Conversion disorder (functional neurological symptom disorder), Todd paralysis (post-ictal), Cerebral palsy (congenital), Hereditary spastic paraplegia
\end{itemize}

\section*{When to See a Doctor}
See your doctor if:
\begin{itemize}
  \item Sudden onset of unilateral weakness (FAST --- Face, Arm, Speech, Time) --- immediate emergency evaluation for acute stroke
  \item Acute or subacute paralysis ascending from legs upward (Guillain-Barr\'e suspicion)
  \item Paralysis after trauma (spinal cord injury or intracranial hemorrhage)
  \item Weakness that fluctuates or worsens with use (myasthenia gravis)
  \item Paralysis accompanied by bowel or bladder dysfunction (spinal cord syndrome)
  \item New-onset paralysis in a child (acute flaccid myelitis)
\end{itemize}

\section*{Related Symptoms}
\begin{itemize}
  \item Sensory loss or paresthesia
  \item Spasticity or flaccidity
  \item Hyperreflexia or hyporeflexia
  \item Babinski sign (extensor plantar response)
  \item Muscle atrophy
  \item Fasciculations
  \item Dysarthria or aphasia
  \item Dysphagia
  \item Bowel/bladder dysfunction
  \item Loss of proprioception
\end{itemize}
\textbf{Important:} Acute paralysis is a medical emergency requiring immediate neuroimaging (non-contrast CT head for stroke); the distinction between UMN and LMN patterns localizes the lesion and guides further workup.

\section*{Self-Care / Home Management}
\begin{itemize}
  \item Acute paralysis is not manageable at home --- seek emergency medical attention immediately.
  \item Chronic paralysis management: prevent skin breakdown (turn/reposition q2h for bedbound people), range of motion exercises to prevent contractures, deep vein thrombosis prophylaxis (Lovenox or compression devices), and bowel/bladder program.
  \item Durable medical equipment: wheelchair, transfer board, hospital bed, commode as needed.
  \item Fall prevention with grab bars, shower chair, and clear pathways.
\end{itemize}

\section*{How Your Doctor Will Evaluate You}
\begin{itemize}
  \item History establishes onset (sudden vs. gradual), progression, pattern (mono/hemi/para/quadriplegia), associated symptoms (pain, sensory loss, bowel/bladder), and risk factors.
  \item Neurological examination differentiates UMN (spasticity, hyperreflexia, Babinski, clonus) from LMN (flaccidity, hyporeflexia, fasciculations, atrophy) and localizes the lesion.
  \item Urgent non-contrast CT head for acute hemiparesis to rule out hemorrhage; CT angiography and perfusion imaging for ischemic stroke.
  \item MRI brain and/or spine with and without contrast for structural, demyelinating, or inflammatory lesions.
  \item Electromyography (EMG) and nerve conduction studies (NCS) for peripheral nerve, neuromuscular junction, and muscle disease.
  \item Lumbar puncture for suspected inflammatory or infectious myelitis (CSF analysis, oligoclonal bands, IgG index).
  \item Serology: AChR antibodies (myasthenia gravis), ganglioside antibodies (Guillain-Barr\'e), GM1 antibodies (multifocal motor neuropathy), viral PCR (West Nile, enterovirus).
\end{itemize}

\section*{Treatment Options}
\begin{itemize}
  \item Ischemic stroke: IV thrombolysis within 4.5 hours, mechanical thrombectomy for large vessel occlusion within 6--24 hours (based on perfusion imaging), antiplatelet therapy, and risk factor management.
  \item Intracerebral hemorrhage: blood pressure control, reversal of anticoagulation, surgical evacuation if indicated.
  \item Guillain-Barr\'e syndrome: IVIG or plasma exchange, respiratory monitoring (PEF, vital capacity), and ICU admission for declining function.
  \item Myasthenia gravis: acetylcholinesterase inhibitors (pyridostigmine), immunosuppression (prednisone, mycophenolate), thymectomy for thymoma; IVIG or plasma exchange for crisis.
  \item Multiple sclerosis relapse: high-dose IV methylprednisolone for acute exacerbations; disease-modifying therapies for prevention.
  \item Spinal cord compression: urgent neurosurgical decompression for traumatic or neoplastic causes; high-dose corticosteroids for malignant cord compression.
  \item Stroke rehabilitation: physical therapy, occupational therapy, speech therapy, constraint-induced movement therapy, and spasticity management (baclofen, botulinum toxin).
  \item Assistive devices and long-term care planning for permanent paralysis.
\end{itemize}

\section*{References}
\begin{thebibliography}{99}
\bibitem{paralysis:paralysis1} Caplan LR. "Acute stroke." In: Caplan\'s Stroke: A Clinical Approach. 5th ed. Elsevier; 2016.
\bibitem{paralysis:paralysis2} Powers WJ, Rabinstein AA, Ackerson T, et al. "Guidelines for the early management of patients with acute ischemic stroke." Stroke. 2019;50(12):e344-e418.
\bibitem{paralysis:paralysis3} GBD 2016 Stroke Collaborators. "Global, regional, and national burden of stroke, 1990-2016." Lancet Neurol. 2019;18(5):439-458.
\bibitem{paralysis:paralysis4} Peacock WF, Rafique Z, Singer AJ. "Acute stroke." In: Rosen\'s Emergency Medicine. 10th ed. Elsevier; 2023.
\bibitem{paralysis:paralysis5} Adams HP, del Zoppo G, Alberts MJ, et al. "Guidelines for the early management of adults with ischemic stroke." Stroke. 2007;38(5):1655-1711.
\bibitem{paralysis:paralysis6} Jauch EC, Saver JL, Adams HP, et al. "Guidelines for the early management of patients with acute ischemic stroke." Stroke. 2013;44(3):870-947.
\end{thebibliography}
